\subsection*{Assessment Overview}

The assessment began with the definition of the testing scope,
objectives, rules of engagement, and target environment. Reconnaissance
activities were then performed to identify application functionality,
endpoints, parameters, and other accessible components.

Automated security testing was performed using tools such as OWASP
ZAP and FFUF, while SQLMap was used where appropriate for SQL
Injection assessment. The identified potential vulnerabilities were
then manually verified using techniques including crafted HTTP
requests and Burp Suite.

Only vulnerabilities that were successfully validated and supported
by appropriate evidence were included as confirmed findings in this
report.

\subsection*{Confirmed Findings}

A total of eight vulnerabilities were confirmed during the
assessment.

\begin{table}[H]
\centering
\renewcommand{\arraystretch}{1.25}
\begin{tabular}{|l|p{6cm}|l|}
\hline
\textbf{Finding ID} & \textbf{Vulnerability} & \textbf{Severity} \\
\hline
BAC & Broken Access Control -- Privilege Escalation via Registration
& Critical \\
\hline
IDOR & Insecure Direct Object Reference -- Unauthorized Basket Access
& Medium \\
\hline
AUTH & Weak Password Recovery & Critical \\
\hline
SQLI-01 & SQL Injection -- Authentication Bypass & Critical \\
\hline
SQLI-02 & SQL Injection -- Product Search and Data Extraction
& High \\
\hline
XSS & Reflected Cross-Site Scripting -- Search Parameter & Medium \\
\hline
CSRF & Cross-Site Request Forgery -- Profile Update & Medium \\
\hline
OR & Open Redirect -- Unvalidated Redirect Parameter & Medium \\
\hline
\end{tabular}
\caption{Summary of confirmed vulnerabilities}
\end{table}

\subsection*{Severity Distribution}

The confirmed findings consist of three Critical-severity findings,
one High-severity finding, and four Medium-severity findings.

\begin{itemize}
    \item \textbf{Critical:} 3 findings
    \item \textbf{High:} 1 finding
    \item \textbf{Medium:} 4 findings
    \item \textbf{Low:} 0 findings
\end{itemize}

The Critical findings include Broken Access Control, SQL Injection --
Authentication Bypass, and Weak Password Recovery. The single
High-severity finding is SQL Injection -- Product Search and Data
Extraction. Reflected XSS, IDOR, CSRF, and Open Redirect were
classified as Medium severity based on the demonstrated impact and
exploitation conditions.

\subsection*{Key Security Observations}

The assessment identified weaknesses across several important
application security controls.

The SQL Injection findings demonstrate insufficient protection of
user-controlled input used by database-related functionality. The
login functionality was found to be susceptible to authentication
bypass, while the product search functionality was found to allow
unauthorized database interaction under the tested conditions.

The Broken Access Control finding demonstrates a weakness in
server-side authorization and privilege assignment. The Weak Password
Recovery finding demonstrates that an attacker can take over the
affected account without knowledge of the original password. If the
same weakness affected a privileged account, it could potentially lead
to unauthorized administrative access.

The remaining findings affect resource authorization, client-side
script
execution, state-changing requests, and redirect handling.
Together, these findings demonstrate that security controls should be
strengthened across multiple areas of the application.

\subsection*{Business and Security Impact}

The confirmed vulnerabilities may affect the confidentiality,
integrity, and security of application resources depending on the
specific vulnerability and exploitation conditions.

The Critical and High severity findings require particular attention
because they affect authentication, authorization, account recovery,
and database interaction.

If similar weaknesses existed in a production application handling
real user information, successful exploitation could result in
unauthorized access, modification of application resources, exposure
of data, account compromise, or execution of attacker-controlled
content.

\subsection*{Recommended Actions}

The identified vulnerabilities should be addressed according to their
severity and associated risk.

Priority should be given to implementing parameterized database
queries, enforcing server-side authorization and role assignment,
strengthening password recovery controls, implementing object-level
authorization checks, and applying appropriate output encoding and
request protection.

Additional controls should include secure redirect validation,
appropriate security headers, least-privilege access, server-side
input validation, and security monitoring.

After remediation, all confirmed findings should be retested using
the original proof-of-concept techniques and additional relevant test
cases to verify that the vulnerabilities have been effectively
addressed.

\subsection*{Report Scope}

This report documents the testing activities, identified
vulnerabilities, supporting evidence, impact assessment, and
recommended remediation measures within the defined assessment scope.

The findings and conclusions presented in this report are based on
the testing performed against the specified OWASP Juice Shop
environment during this engagement.